\documentclass[11pt,reqno]{amsart}
\usepackage{amsmath,amssymb,graphicx,amscd}
\usepackage{mathrsfs}
\usepackage{enumerate}
% pour les S des spheres
\usepackage{mathrsfs}
\usepackage{dsfont}
\usepackage{stmaryrd}
\usepackage{wrapfig}
\usepackage{pstricks}

% ----------------------------------------------------------------
\vfuzz2pt % Don't report over-full v-boxes if over-edge is small
\hfuzz2pt % Don't report over-full h-boxes if over-edge is small

\newcommand{\prob}{{\mathbb{P}}}

\renewcommand{\S}{\mathscr{S}}  % complex unit sphere

\newcommand{\CC}{{\mathbb C}}
\newcommand{\HH}{{\mathbb H}}
\newcommand{\RR}{{\mathbb R}}
\newcommand{\QQ}{{\mathbb{Q}}}
\newcommand{\NN}{{\mathbb{N}}}
\newcommand{\ZZ}{{\mathbb{Z}}}
\renewcommand{\P}{{\mathbb{P}}}
\renewcommand{\span}{{\mathrm{span}}}
\newcommand{\sgn}{\operatorname{{sgn}}}
 \DeclareMathOperator{\Prob}{\mathbb{P}}   %probability
 \DeclareMathOperator{\E}{\mathbb{E}}      %expectation
 \DeclareMathOperator{\cov}{cov}             %logarithmic integral

 \newcommand{\I}{1\!\!1}                   %Indicator function


 \newcommand{\dd}{{\mathrm{d}}}            %differential
 \newcommand{\ii}{{\mathrm{i}}}
  \newcommand{\Id}{{\mathrm{Id}}}
\newcommand{\law}{\overset{\mbox{\rm \scriptsize law}}{=}}
\newcommand{\convlaw}{\overset{\mbox{\rm \scriptsize law}}{\longrightarrow}}


\renewcommand{\Re}{{\mathfrak{Re}}}
\renewcommand{\Im}{{\mathfrak{Im}}}

% THEOREMS -------------------------------------------------------

\newtheorem{thm}{Theorem}[section]
\newtheorem{cor}[thm]{Corollary}
\newtheorem{lem}[thm]{Lemma}
\newtheorem{prop}[thm]{Proposition}
\theoremstyle{definition}
\newtheorem{defn}[thm]{Definition}
\theoremstyle{remark}
\newtheorem*{rem}{Remark}



\numberwithin{equation}{section}



\begin{document}
\begin{center}
{\bf Probability, homework 1, due September 9.}
\end{center}
\vspace{2cm}
\small

\noindent{\bf Exercise 1}.
Let $(\mathcal{G}_\alpha)_{\alpha\in A}$ be an arbitrary family of $\sigma$-fields defined on an abstract space $\Omega$, with $A$ possibly uncountable.  Show that $\cap_{\alpha\in A}\mathcal{G}_\alpha$ is also a $\sigma$-field.\\





\vspace{0.2cm}

\noindent{\bf Exercise 2}. Let $\varnothing\subsetneq A\subsetneq B\subsetneq \Omega$ (these are strict inclusions). What is the $\sigma$-field generated by $\{A,B\}$ ?\\


\vspace{0.2cm}

\noindent{\bf Exercise 3}.  Let $\mathcal{F}$, $\mathcal{G}$ be $\sigma$-fields for the same $\Omega$.  Is $\mathcal{F}\cup\mathcal{G}$ a $\sigma$-field?\\

\vspace{0.2cm}

\noindent{\bf Exercise 4}.  For $\Omega=\mathbb{N}$ and $n\geq 0$, let $\mathcal{F}_n=\sigma(\{\{0\},\dots,\{n\}\})$. Show that $(\mathcal{F}_n)_{n \geq 0}$ is a non-decreasing sequence but that $\cup_{n\geq 0}\mathcal{F}_n$ is a not a $\sigma$-field.\\





\noindent{\bf Exercise 5}. 
Let $\Omega$ be an infinite set (countable or not). Let $\mathcal{A}$ be the set of subsets of $\Omega$ that are either finite
or with finite complement in $\Omega$. Prove that $\mathcal{A}$ is a field but not a $\sigma$-field.\\


\noindent{\bf Exercise 6}. 
Can you build an infinite,  countable $\sigma$-field?\\

\noindent{\bf Exercise 7}. 
A monotone class is a collection $\mathcal{M}$ of sets closed under both monotone increasing and monotone decreasing (i.e.  if $A_i\in\mathcal{M}$ and either $A_i\uparrow A$ or $A_i\downarrow A$ then $A\in\mathcal{M}$)

Prove that if $\mathcal{A}\subset \mathcal{M}$ with $\mathcal{A}$ a field and $\mathcal{M}$ a monotone class,  then $\sigma(\mathcal{A})\subset \mathcal{M}$.\\


\noindent{\bf Exercise 8}. Let $\mathbb{P}$ be a probability measure on $\Omega$, endowed with a $\sigma$-field $\mathcal{A}$. 
\begin{enumerate}[(i)]
\item What is the meaning of the following events, where all $A_n$'s are elements of $\mathcal{A}$?
$$
\liminf_{n\to\infty}A_n=\bigcup_{n\geq 1}\bigcap_{k\geq n} A_k,\ \
\limsup_{n\to\infty}A_n=\bigcap_{n\geq 1}\bigcup_{k\geq n} A_k.
$$
\item Prove that $\limsup_{n\to\infty} A_n$ and $\liminf_{n\to\infty} A_n$ are in $\mathcal{A}$.
\item In the special case $\Omega=\RR$, for any $p\geq 1$, let
$$
A_{2p}=\left[-1,2+\frac{1}{2p}\right),\ \ A_{2p+1}=\left(-2-\frac{1}{2p+1},1\right].
$$
What are $\liminf_{n\to\infty}A_n$ and $\limsup_{n\to\infty}A_n$?
\item Prove that the following always holds:
$$
\mathbb{P}\left(\liminf_{n\to\infty} A_n\right)\leq \liminf_{n\to\infty}\mathbb{P}\left(A_n\right),
\mathbb{P}\left(\limsup_{n\to\infty} A_n\right)\geq \limsup_{n\to\infty}\mathbb{P}\left(A_n\right).
$$
\end{enumerate}
\vspace{0.2cm}


\noindent{\bf Exercise 9}. The symmetric difference of two events $A$ and $B$, denoted $A\triangle B$, is the event that precisely one of them occurs: $A\triangle B=(A\cup B)\backslash(A\cap B)$.
\begin{enumerate}[(i)]
\item Write a formula for $A\triangle B$ that only involves the operations of union, intersection and complement, but no set difference.
\item Define $d(A,B)=\P(A\triangle B)$. Show that for any three events $A$, $B$, $C$, 
$$
d(A,B)+d(B,C)-d(A,C)
=
2\left(\P\left(A\cap B^c\cap C\right)+\P\left(A^c\cap B\cap C^c\right)\right).
$$
\item Assume $A\subset B\subset C$. Prove that $d(A,C)=d(A,B)+d(B,C)$.\\
\end{enumerate}

\noindent{\bf Exercise 10}.
Prove the Bonferroni inequalities: if $A_i\in \mathcal{A}$ is a sequence of events, then
\begin{enumerate}[(i)]
\item $\P\left(\cup_{i=1}^n A_i\right)\geq \sum_{i=1}^n\P(A_i)-\sum_{i<j}\P(A_i\cap A_j)$,
\item $\P\left(\cup_{i=1}^n A_i\right)\leq \sum_{i=1}^n\P(A_i)-\sum_{i<j}\P(A_i\cap A_j)
+\sum_{i<j<k}\P(A_i\cap A_j\cap A_k)$.
\end{enumerate} 


\end{document} 